\section{基于Boosting的集成}
\label{sec:ensemble-boosting}

Bagging中的成员可以同时训练，因为一个成员不需要知道其他成员的预测。若希望新增成员专门补充当前组合尚未处理好的部分，这种独立性就不再成立。\term{Boosting}（提升）构造一个序贯加法模型：第$t$轮先观察已有函数$F_{t-1}$，再选择新成员及其贡献，使指定训练目标继续下降。AdaBoost、梯度提升及现代提升树共享这种依赖关系，但它们生成修正信号和近似目标的方式不同。

\subsection{前向逐步加法模型}
\label{sec:ensemble-stagewise-additive}

设$F_0$为初始预测，Boosting经过$T$轮得到
\begin{equation}
  F_T(\bm{x})
  =
  F_0(\bm{x})
  +
  \sum_{t=1}^{T}\alpha_t h_t(\bm{x}),
  \label{eq:ensemble-additive-model}
\end{equation}
其中$h_t$是第$t$个成员，$\alpha_t$控制其贡献。为使复杂度惩罚对应真正加入模型的函数，把本轮增量记为$f_t=\alpha_t h_t$。给定损失$\ell$，理想化的第$t$轮直接优化增量：
\begin{equation}
  f_t
  \in
  \argmin_{f\in\mathcal{F}}
  \sum_{i=1}^{n}
  \ell\!\left(
    y_i,
    F_{t-1}(\bm{x}_i)+f(\bm{x}_i)
  \right)
  +\Omega_t(f).
  \label{eq:ensemble-forward-stagewise}
\end{equation}
这里$\mathcal{F}=\{\alpha h:\alpha\in\mathbb{R},h\in\mathcal{H}\}$是由成员函数族$\mathcal{H}$生成的增量函数族，$\Omega_t$是相应复杂度惩罚，也可以取零。对于树，叶数惩罚限制结构规模，叶输出的范数惩罚限制增量幅度；后者必须作用于$f$的叶输出，不能只惩罚仍可与$\alpha$相互缩放的$h$。前向逐步训练固定已经加入的成员，只优化本轮增量；它避免一次性联合搜索全部函数，却也不等于求得式\eqref{eq:ensemble-additive-model}在所有成员和权重上的全局最优解。后文的梯度拟合、二阶近似和额外收缩也不等于精确求解这一理想目标。

不同Boosting方法的主要差别，可以从“当前不足怎样变成下一轮训练信号”来判断。AdaBoost改变训练样本的权重，使当前错分记录在下一轮具有更大影响；梯度提升把损失对当前预测的负梯度作为伪响应；二阶提升树还利用局部曲率近似叶权重和分裂收益。只有平方损失下的负梯度等于普通响应残差，把所有Boosting概括为“反复拟合残差”会掩盖损失函数的作用。

\subsection{AdaBoost}
\label{sec:ensemble-adaboost}

\term{AdaBoost}（adaptive boosting）从带权训练样本出发，使每一轮弱分类器回应当前加权分布。考虑二分类标签$y_i\in\{-1,1\}$和成员$h_t(\bm{x})\in\{-1,1\}$。初始样本权重为$w_i^{(1)}=1/n$，第$t$轮成员的加权错误率为
\begin{equation}
  \epsilon_t
  =
  \sum_{i=1}^{n}
  w_i^{(t)}
  \mathbb{I}\!\left\{y_i\neq h_t(\bm{x}_i)\right\}.
  \label{eq:ensemble-adaboost-error}
\end{equation}
当$0<\epsilon_t<1/2$时，成员投票权重取
\begin{equation}
  \alpha_t
  =
  \frac{1}{2}
  \log\frac{1-\epsilon_t}{\epsilon_t},
  \qquad
  w_i^{(t+1)}
  =
  \frac{
    w_i^{(t)}
    \exp\!\left[-\alpha_t y_i h_t(\bm{x}_i)\right]
  }{Z_t},
  \label{eq:ensemble-adaboost-update}
\end{equation}
其中$Z_t$使新权重和为1。错分样本满足$y_i h_t(\bm{x}_i)=-1$，权重乘上$e^{\alpha_t}$；正确样本的权重乘上$e^{-\alpha_t}$。若$\epsilon_t=0$，在这些正权重下当前成员已分对全部训练样本，便直接输出$h_t$并停止，不再计算发散的$\alpha_t$或执行归一化更新；这与\BookSectionRef{sec:boosting}{第二卷“基础、理论和可信性”的“二分类学习的可学习性”章“集成学习”节}的约定一致。其余情况下，最终分类器在预先固定平局规则后按$\operatorname{sign}(\sum_t\alpha_t h_t(\bm{x}))$决策。若成员不能在当前权重下取得小于$1/2$的错误率，标准更新就无法获得正的成员贡献，需要停止、调整成员训练或采用相应变体。

图\ref{fig:ensemble-adaboost-weight-update}展示了式\eqref{eq:ensemble-adaboost-update}形成的反馈回路。当前分类结果先改变样本权重，再由归一化后的权重决定下一轮成员看到的训练分布；错分记录因而获得更大的后续影响。

\begin{figure}[htbp]
  \centering
  % AdaBoost turns classification outcomes into the next round's sample weights.
% Schematic only; no empirical quantities are encoded.
% Canvas: 112 mm x 61 mm. Dependencies: project palette and TikZ.
\begin{tikzpicture}[
  x=1mm,y=1mm,
  font=\figurefont\footnotesize,
  text=FigureInk,
  stage/.style={
    rectangle,
    model hidden,
    minimum width=24mm,
    minimum height=11mm,
    text width=21mm,
    align=center,
    inner sep=1mm
  },
  data/.style={
    stage,
    model input
  },
  model/.style={
    stage,
    model hidden
  },
  update/.style={
    rectangle,
    model training,
    minimum width=29mm,
    minimum height=11mm,
    text width=26mm,
    align=center,
    inner sep=1mm
  },
  flow/.style={
    model flow
  },
  feedback/.style={
    model feedback
  }
]
  \path[use as bounding box] (0,0) rectangle (112,61);

  \node[anchor=west,font=\figurefont\small] at (1,57)
    {第$t$轮：分类结果决定第$t+1$轮的训练分布};

  \node[data] (weighted-data) at (15,40)
    {带权训练样本\\$\{w_i^{(t)}\}$};
  \node[model] (learner) at (43,40)
    {训练弱分类器\\$h_t$};
  \node[update] (outcome) at (88,49)
    {正确分类\\乘以$e^{-\alpha_t}$};
  \node[update] (mistake) at (88,31)
    {错分\\乘以$e^{\alpha_t}$};
  \node[update,minimum width=29mm] (normalize) at (44,12)
    {归一化\\得到$\{w_i^{(t+1)}\}$};

  \draw[flow] (weighted-data)--(learner);
  \draw[FigureInk,line width=.9pt] (learner.east)--(63,40);
  \fill[FigureInk] (63,40) circle[radius=1.1pt];
  \draw[flow] (63,40)|-(outcome.west);
  \draw[feedback] (63,40)|-(mistake.west);

  \draw[FigureInk,line width=.9pt]
    (outcome.east)--(109,49)--(109,20)--(106,20);
  \draw[FigureRed,dashed,line width=.9pt]
    (mistake.east)--(106,31)--(106,20);
  \fill[FigureInk] (106,20) circle[radius=1.1pt];
  \draw[flow] (106,20)--(65,20)--(65,12)--(normalize.east);
  \draw[feedback] (normalize.west)--(1,12)|-(weighted-data.west);
\end{tikzpicture}

  \caption{AdaBoost的样本权重更新。分类结果把样本分成正确与错分两路，两路缩放后的权重汇合并归一化，形成下一轮训练分布。实线表示常规训练与正确分类样本的更新路径，红色虚线强调错分样本被放大的路径。}
  \label{fig:ensemble-adaboost-weight-update}
\end{figure}

这一更新也可以从指数损失解释。令$F(\bm{x})=\sum_t\alpha_t h_t(\bm{x})$，经验目标为
\[
  \sum_{i=1}^{n}\exp\!\left[-y_iF(\bm{x}_i)\right].
\]
固定$F_{t-1}$后选择$h_t$与$\alpha_t$，在$0<\epsilon_t<1/2$时可以导出式\eqref{eq:ensemble-adaboost-update}。因此，样本重加权不是与目标函数无关的经验技巧，而是指数损失下前向逐步加法建模的结果\scite{freund1997boosting,friedman2000adaboost}

指数权重会持续放大难以正确分类的训练记录。若这些记录代表真实而稀少的边界情形，重新分配注意力可能有用；若它们来自标签错误或异常观测，后续成员可能反复追逐不可修正的噪声。训练错误下降还只说明组合在当前样本上的拟合变化，不能替代新样本风险。弱学习条件、训练错误界及投票间隔的泛化分析见
\BookSectionRef{sec:boosting}{第二卷“基础、理论和可信性”的“二分类学习的可学习性”章“集成学习”节}；
这里更关心更新如何形成可使用的模型。

多分类AdaBoost需要重新定义成员错误与类别得分，回归变体则要用响应误差决定样本权重和预测聚合。它们保留序贯重加权思想，却不再直接使用上述二分类推导。比较这些变体时，应按第\ref{chap:classification}章或第\ref{chap:regression}章的输出和损失分别评价，不能只沿用二分类准确率。

\subsection{梯度提升机与GBDT}
\label{sec:ensemble-gradient-boosting}

AdaBoost的指数损失给出了特定重加权规则。若任务需要平方损失、对数损失或其他关于预测可微的目标，可以直接询问：当前预测沿哪个函数方向变化，经验损失下降最快？\term{梯度提升机}（gradient boosting machine, GBM）用基学习器逼近这个负梯度方向，从而把序贯加法扩展到更一般的损失\scite{friedman2001gradientboost}

通常从使训练损失最小的常数预测开始：
\begin{equation}
  F_0(\bm{x})=c_0,
  \qquad
  c_0\in\argmin_c\sum_{i=1}^{n}\ell(y_i,c).
  \label{eq:ensemble-gbm-initialization}
\end{equation}
平方损失下$c_0$是响应均值。这里要求存在有限的最小化点；若某种损失与数据组合使最优常数趋于无穷，就需约束预测范围或另定有限初始化。

在第$t$轮，对每个训练样本计算\term{伪残差}（pseudo-residual），即损失在当前预测处的负导数：
\begin{equation}
  r_i^{(t)}
  =
  -
  \left.
  \frac{\partial \ell(y_i,z)}{\partial z}
  \right|_{z=F_{t-1}(\bm{x}_i)}.
  \label{eq:ensemble-pseudo-residual}
\end{equation}
随后用平方误差把这些修正量拟合到成员函数族：
\begin{equation}
  h_t\in\argmin_{h\in\mathcal{H}}
  \sum_{i=1}^{n}\bigl(r_i^{(t)}-h(\bm{x}_i)\bigr)^2.
  \label{eq:ensemble-gbm-gradient-fit}
\end{equation}
这是对负梯度的平方拟合，不是用原任务损失把$r_i^{(t)}$当成新标签；树的贪心生长通常只近似求解这一子问题。固定$h_t$后，可以按原损失做整体线搜索，在最小值存在时取
\begin{equation}
  \alpha_t\in\argmin_\alpha
  \sum_{i=1}^{n}
  \ell\!\left(y_i,F_{t-1}(\bm{x}_i)+\alpha h_t(\bm{x}_i)\right).
  \label{eq:ensemble-gbm-line-search}
\end{equation}
此处$\alpha_t$是整棵树或整个成员共享的一个标量。加入预先选择的收缩系数$\eta\in(0,1]$后，
\begin{equation}
  F_t(\bm{x})
  =
  F_{t-1}(\bm{x})
  +
  \eta\alpha_t h_t(\bm{x}).
  \label{eq:ensemble-gradient-boosting-update}
\end{equation}
当$\ell(y,z)=\tfrac12(y-z)^2$时，式\eqref{eq:ensemble-pseudo-residual}给出$y_i-F_{t-1}(\bm{x}_i)$，恰好是普通残差。分类时则要先明确得分与概率的关系。例如二分类采用$y_i\in\{0,1\}$，令$F$表示对数几率（logit），相应概率与二项对数损失为
\[
  p=\sigma(F)=\frac{1}{1+e^{-F}},\qquad
  \ell(y,F)=\log(1+e^F)-yF.
\]
由于$\partial\ell/\partial F=\sigma(F)-y$，第$t$轮的伪残差为
\[
  r_i^{(t)}=y_i-p_i^{(t-1)},\qquad
  p_i^{(t-1)}=\sigma\!\left(F_{t-1}(\bm{x}_i)\right).
\]
例如$y_i=1$而当前概率为$0.3$时，伪残差是$0.7$。成员拟合的是损失关于得分的连续修正量，增量按式\eqref{eq:ensemble-gradient-boosting-update}加到$F$上，再由Sigmoid换算为新概率；不能把$0.7$直接加到概率上。

若$h_t$采用回归树，就得到\term{梯度提升决策树}（gradient boosting decision tree, GBDT）。除了整体线搜索，还可以保留平方拟合得到的叶区域$\bm{R}_1^{(t)},\ldots,\bm{R}_{L_t}^{(t)}$，再按原损失分别求每片叶的增量。当各叶内存在有限最小化点时，
\begin{equation}
  \begin{aligned}
    v_j^{(t)}
      &\in\argmin_v
        \sum_{i:\,\bm{x}_i\in\bm{R}_j^{(t)}}
        \ell\!\left(y_i,F_{t-1}(\bm{x}_i)+v\right),\\
    F_t(\bm{x})
      &=F_{t-1}(\bm{x})
        +\eta\sum_{j=1}^{L_t}
        v_j^{(t)}\mathbb{I}\{\bm{x}\in\bm{R}_j^{(t)}\}.
  \end{aligned}
  \label{eq:ensemble-gbm-leaf-update}
\end{equation}
这里$L_t$是叶数，各叶的$v_j^{(t)}$替换原树的输出，不是同一个$\alpha_t$的不同写法；这种逐叶更新替代整体线搜索，而非在其后再乘一轮叶权重。若叶内最优值只能在无穷处逼近，则需限制叶输出或引入相应正则，不能直接把无穷值代入更新。“回归树”描述成员拟合连续修正量的计算角色，整体任务仍可以是分类。

第\ref{sec:regression-quantiles}节的分位数损失$\ell(y,z)=\rho_\tau(y-z)$在$z=y$处不可微，不能处处使用式\eqref{eq:ensemble-pseudo-residual}。分位数提升可改用负次梯度，例如约定
\begin{equation}
  r_i^{(t)}
  =\tau-\mathbb{I}\{y_i<F_{t-1}(\bm{x}_i)\},
  \qquad 0<\tau<1.
  \label{eq:ensemble-quantile-pseudo-residual}
\end{equation}
等号处的负次梯度可取$[\tau-1,\tau]$内任意值，这里选端点$\tau$。常数初始化取样本$\tau$分位数，仍用平方误差拟合上述修正量；逐叶更新则最小化原分位数损失，故$v_j^{(t)}$取该叶内普通残差$y_i-F_{t-1}(\bm{x}_i)$的$\tau$分位数，而不是伪残差的均值。

学习率、轮数与单树复杂度共同决定加法函数。较小学习率缩小每轮变化，通常需要更多轮才能达到相近训练目标；较深的树可以在单轮表达更高阶交互，也会扩大每轮容量。行采样和列采样进一步引入随机性，早停则依据独立验证轨迹选择轮数。不存在只调其中一个参数就能固定其他参数作用的普遍规则。

GBM的轮次具有数据依赖，不能像Bagging成员那样全部并行。单轮内部仍可并行统计不同特征或候选分裂。训练损失随轮次下降、优化过程趋于稳定和测试风险达到较小值是三件不同的事；若轮数、损失、树结构和早停规则都根据同一验证集反复调整，还要把这种选择纳入最终评价。

\subsection{XGBoost}
\label{sec:ensemble-xgboost}

经典梯度提升用一阶负梯度给出拟合方向。\term{XGBoost}在正则化加法树目标中进一步使用二阶局部近似，并围绕稀疏输入、近似分裂和内存访问设计训练系统\scite{chen2016xgboost}这三层需要分开理解：二阶目标决定怎样评价一棵候选树，分裂算法近似搜索树结构，系统机制则决定统计量怎样高效计算。

设第$t$棵树为$f_t$，损失在当前预测$F_{t-1}(\bm{x}_i)$处的一阶导与二阶导分别为$g_i^{(t)}$和$s_i^{(t)}$。去掉与$f_t$无关的常数后，本轮目标近似为
\begin{equation}
  \widetilde{\mathcal{L}}^{(t)}
  =
  \sum_{i=1}^{n}
  \left[
    g_i^{(t)}f_t(\bm{x}_i)
    +\frac{1}{2}s_i^{(t)}f_t^2(\bm{x}_i)
  \right]
  +\gamma L
  +\frac{\lambda}{2}\sum_{j=1}^{L}v_j^2,
  \label{eq:ensemble-xgboost-objective}
\end{equation}
其中$L$是叶数，$v_j$是第$j$片叶的输出，$\gamma$惩罚新增叶，$\lambda$约束叶权重。若$I_j$表示落入第$j$片叶的样本索引，记
\[
  G_j=\sum_{i\in I_j}g_i^{(t)},
  \qquad
  S_j=\sum_{i\in I_j}s_i^{(t)}.
\]
在$S_j+\lambda>0$时，固定树结构的最优叶权重为
\begin{equation}
  v_j^\star=-\frac{G_j}{S_j+\lambda}.
  \label{eq:ensemble-xgboost-leaf-weight}
\end{equation}
把一个父叶分成左右子叶时，相应近似目标下降量为
\begin{equation}
  \operatorname{Gain}
  =
  \frac{1}{2}
  \left[
    \frac{G_{\mathrm{L}}^2}{S_{\mathrm{L}}+\lambda}
    +
    \frac{G_{\mathrm{R}}^2}{S_{\mathrm{R}}+\lambda}
    -
    \frac{(G_{\mathrm{L}}+G_{\mathrm{R}})^2}
         {S_{\mathrm{L}}+S_{\mathrm{R}}+\lambda}
  \right]
  -\gamma.
  \label{eq:ensemble-xgboost-split-gain}
\end{equation}
式\eqref{eq:ensemble-xgboost-split-gain}把当前损失的梯度、曲率与树复杂度放在同一个分裂准则中。它是局部二阶近似下的增益，不表示贪心生长已经找到所有树结构中的全局最优。

图\ref{fig:ensemble-gradient-xgboost}把一阶梯度提升与XGBoost的二阶树评分放在同一计算链中。左图回答新成员拟合什么，右图回答候选树结构怎样利用梯度统计得到叶权重和分裂增益；二阶统计扩展了本轮目标的近似信息，但没有改变逐轮加法更新的外层结构。

\begin{figure*}[htbp]
  \centering
  % Gradient boosting's first-order signal and XGBoost's second-order tree score.
% Schematic only; no empirical quantities are encoded.
% Canvas: 169 mm x 64 mm. Dependencies: project palette and TikZ.
\begin{tikzpicture}[
  foundation picture,x=1mm,y=1mm,
  font=\figurefont\footnotesize,
  state/.style={
    foundation node,
    draw=FigureBlue,
    fill=FigureBlueFill,
    minimum width=26mm,
    minimum height=10mm,
    text width=22mm,
    align=center,
    inner sep=1mm
  },
  model/.style={
    foundation node,
    draw=FigureRed,
    fill=FigureRedFill,
    minimum width=28mm,
    minimum height=10mm,
    text width=24mm,
    align=center,
    inner sep=1mm
  },
  signal/.style={
    foundation node,
    draw=FigureRed,
    fill=white,
    minimum width=30mm,
    minimum height=11mm,
    text width=26mm,
    align=center,
    inner sep=1mm
  },
  result/.style={
    foundation node,
    draw=FigureYellow,
    fill=FigureYellowFill,
    minimum width=29mm,
    minimum height=11mm,
    text width=25mm,
    align=center,
    inner sep=1mm
  },
  flow/.style={
    foundation flow
  }
]
  \path[use as bounding box] (0,0) rectangle (169,64);

  \node[anchor=west,foundation heading] at (1,60)
    {(a) 梯度提升：把损失方向变成新成员};
  \node[state] (current) at (18,42)
    {当前预测\\$F_{t-1}(\bm{x}_i)$};
  \node[signal] (gradient) at (52,42)
    {计算负梯度\\$r_i^{(t)}=-\partial_z\ell$};
  \node[model] (fit) at (52,19)
    {用$h_t$拟合\\伪残差};
  \node[result] (updated) at (18,19)
    {加法更新\\$F_t=F_{t-1}+\eta\alpha_t h_t$};
  \draw[flow] (current)--(gradient);
  \draw[flow] (gradient)--(fit);
  \draw[flow] (fit)--(updated);
  \draw[flow] (updated)--(current);

  \draw[FigureGrid,line width=0.6pt] (80,3)--(80,61);

  \node[anchor=west,foundation heading] at (85,60)
    {(b) XGBoost：二阶统计进入树评分};
  \node[signal] (derivatives) at (97,44)
    {样本统计\\$(g_i^{(t)},s_i^{(t)})$};
  \node[model] (partition) at (147,44)
    {候选分裂\\$I\to(I_{\mathrm L},I_{\mathrm R})$};
  \node[signal] (sums) at (99,22)
    {按候选叶聚合\\$(G_j,S_j)$};
  \node[result,minimum width=22mm,text width=18mm] (leaf) at (145,31)
    {叶权重\\$v_j^\star$};
  \node[result,minimum width=22mm,text width=18mm] (gain) at (145,13)
    {分裂增益\\$\operatorname{Gain}$};

  \draw[flow] (derivatives.east)--(partition.west);
  \draw[flow] (partition.south)--(147,35)--(99,35)--(sums.north);
  \draw[FigureInk,line width=.9pt] (sums.east)--(121,22);
  \fill[FigureInk] (121,22) circle[radius=1.1pt];
  \draw[flow] (121,22)|-(leaf.west);
  \draw[flow] (121,22)|-(gain.west);
\end{tikzpicture}

  \caption{梯度提升与XGBoost中的本轮修正。左图中，损失对当前预测的负梯度形成伪残差，新成员拟合该信号后更新组合；右图中，XGBoost按候选分裂聚合一阶导与二阶导，再由同一组叶统计计算叶权重和分裂增益。复杂度项$\gamma T+\lambda\lVert\bm{w}\rVert_2^2/2$使$\lambda$约束叶权重，并使$\gamma$惩罚新增叶。该图描述局部优化机制，不表示贪心树结构是全局最优。}
  \label{fig:ensemble-gradient-xgboost}
\end{figure*}

精确分裂搜索遍历排序后的候选特征值；近似搜索则先用带权分位数草图构造候选切分，再汇总梯度统计。稀疏感知搜索为缺失值学习默认方向，列采样减少每轮需要检查的特征。块式存储、缓存感知访问、并行统计和外存计算改变的是实现代价，不改变式\eqref{eq:ensemble-xgboost-objective}所定义的学习对象。工程改进能否转化为当前硬件和预算下的速度收益，需要实测。

\subsection{LightGBM}
\label{sec:ensemble-lightgbm}

当样本和特征很多时，逐值搜索分裂会反复扫描大量候选。\term{LightGBM}以直方图训练为基础，把连续特征映射到有限个桶，并在桶级别累计梯度和曲率。构造直方图仍要访问节点样本，但候选分裂只需扫描桶；已知父节点和一个子节点的直方图时，还可以用差分得到另一个子节点的统计量。分桶减少存储和候选数，也引入离散化误差。

树的生长顺序进一步改变计算分配。按层生长同时扩展某一深度上的叶，按叶生长则每次选择当前增益最大的叶继续分裂。在相同叶数下，后者把容量集中到训练目标下降较大的局部区域，可能形成深度不均衡的树。叶数、叶内最小样本量和最大深度因此需要联合约束；训练损失下降更快不等于测试风险必然更小。

LightGBM论文进一步提出两种针对大规模数据的近似。\term{单边梯度采样}（gradient-based one-side sampling, GOSS）保留绝对梯度较大的样本，从绝对梯度较小的样本中随机抽取一部分，并在分裂统计中重新加权后者；它近似的是用于选择分裂的数据统计。例如，在小梯度组内等概率抽取一半样本，就把被抽中者对梯度和的贡献乘二，以补偿未采样部分。一般地，样本以概率$p>0$被保留时，其贡献乘以$1/p$；给定本轮梯度和候选分区，这种逆包含概率加权可无偏估计相应的梯度和，但不意味着由平方、比值或择优操作得到的分裂增益也无偏。

\term{互斥特征捆绑}（exclusive feature bundling, EFB）把很少同时非零的稀疏特征编码到较少的组合特征中，近似减少有效特征数。GOSS、EFB和直方图分别作用于样本、特征与候选阈值，不能合并成一个笼统的“降采样”步骤\scite{ke2017lightgbm}

图\ref{fig:ensemble-lightgbm-approximations}按近似对象并列这三种机制。直方图压缩候选阈值，GOSS压缩参与分裂统计的样本，EFB压缩稀疏特征表示；读者据此可以分别追踪每种近似改变了哪一项计算。

\begin{figure*}[htbp]
  \centering
  % LightGBM reduces split-search work along threshold, sample, and feature axes.
% Schematic only; marks and sparsity patterns encode no measured values.
% Canvas: 164 mm x 64 mm. Dependencies: project palette and TikZ.
\begin{tikzpicture}[
  x=1mm,y=1mm,
  font=\figurefont\footnotesize,
  text=FigureInk,
  flow/.style={
    model flow
  },
  keep/.style={
    circle,
    model training,
    minimum size=5mm,
    inner sep=0pt
  },
  sample/.style={
    circle,
    model input,
    minimum size=4mm,
    inner sep=0pt
  }
]
  \path[use as bounding box] (0,0) rectangle (164,64);

  \node[font=\figurefont\bfseries\small] at (27,59) {(a) 直方图：减少候选阈值};
  \draw[FigureMuted,line width=0.7pt] (5,35)--(49,35);
  \foreach \x in {8,12,15,18,22,24,28,31,35,39,42,46} {
    \fill[NNInput] (\x,35) circle[radius=1.1mm];
  }
  \foreach \x in {5,16,27,38,49} {
    \draw[NNHidden,line width=0.85pt] (\x,28)--(\x,43);
  }
  \node[model output,
    minimum width=38mm,minimum height=9mm,align=center] (bins) at (27,10)
    {桶级梯度统计};
  \draw[flow] (27,22)--(bins);

  \draw[FigureGrid,line width=0.6pt] (54.5,3)--(54.5,61);

  \node[font=\figurefont\bfseries\small] at (82,59) {(b) GOSS：减少参与统计的样本};
  \node[anchor=east,text=FigureMuted] at (72,47) {较大的$|g_i|$};
  \foreach \x in {78,86,94,102} {
    \node[keep] at (\x,47) {};
  }
  \node[anchor=east,text=FigureMuted] at (72,34) {较小的$|g_i|$};
  \foreach \x in {75,81,87,93,99,105} {
    \node[sample] at (\x,34) {};
  }
  \node[model output,
    minimum width=42mm,minimum height=15mm,text width=38mm,align=center]
    (goss) at (82,12)
    {全留大梯度\\抽样并重加权小梯度};
  \draw[flow,draw=NNGradient] (105,47)--(108,47)|-(goss.east);
  \draw[flow] (87,31.5)--(87,27)--(57,27)|-(goss.west);
  \draw[FigureMuted,dashed,line width=0.7pt] (81,31)--(81,37);
  \draw[FigureMuted,dashed,line width=0.7pt] (93,31)--(93,37);
  \draw[FigureMuted,dashed,line width=0.7pt] (105,31)--(105,37);

  \draw[FigureGrid,line width=0.6pt] (109.5,3)--(109.5,61);

  \node[font=\figurefont\bfseries\small] at (137,59) {(c) EFB：减少有效特征};
  \node[anchor=east,text=FigureMuted] at (121,45) {$x_a$};
  \node[anchor=east,text=FigureMuted] at (121,35) {$x_b$};
  \foreach \x in {125,131,137,143,149,155} {
    \draw[FigureGrid,line width=0.6pt] (\x-2.2,42.5) rectangle (\x+2.2,47.5);
    \draw[FigureGrid,line width=0.6pt] (\x-2.2,32.5) rectangle (\x+2.2,37.5);
  }
  \fill[NNInput] (122.8,42.5) rectangle (127.2,47.5);
  \fill[NNInput] (140.8,42.5) rectangle (145.2,47.5);
  \fill[NNInput] (128.8,32.5) rectangle (133.2,37.5);
  \fill[NNInput] (146.8,32.5) rectangle (151.2,37.5);
  \node[model output,
    minimum width=42mm,minimum height=10mm,text width=38mm,align=center]
    (bundle) at (137,10)
    {捆绑特征};
  \draw[flow] (139,21)--(bundle);
\end{tikzpicture}

  \caption{LightGBM中三种作用位置不同的计算近似。直方图把连续取值映射到有限桶，GOSS保留$|g_i|$较大的样本，并抽取、重加权部分$|g_i|$较小的样本，EFB把很少同时非零的特征编码为较少的捆绑特征。图中点数、桶数与非零位置均为机制示意。}
  \label{fig:ensemble-lightgbm-approximations}
\end{figure*}

这些机制的效果依赖梯度分布、特征稀疏性、桶数、冲突处理和硬件实现。论文中的速度与精度比较支持相应实验条件下的结论，不构成LightGBM相对其他实现的无条件排序。公平比较应固定损失、搜索预算和早停规则，并记录离散化与采样设置。

\subsection{CatBoost}
\label{sec:ensemble-catboost}

高基数类别特征若全部展开为独热变量，会增加特征规模；把类别直接映射为整数，又会引入原本不存在的次序。另一种做法是用同类别样本的响应均值编码，但若样本自己的$y_i$参与其编码，模型就能从输入中间接看到待预测标签。这种\term{目标泄漏}（target leakage）会使训练特征与部署特征的构造条件不同。

\term{CatBoost}用随机排列建立\term{有序目标统计}（ordered target statistic）。设$c_i$为样本$i$的某个类别值，用$\pi(j)$表示样本$j$在随机排列中的位置；对当前样本只使用排列中位于它之前且类别相同的记录：
\begin{equation}
  z_i
  =
  \frac{
    \sum_{j:\,\pi(j)<\pi(i)}
    \mathbb{I}\{c_j=c_i\}y_j
    +ap
  }{
    \sum_{j:\,\pi(j)<\pi(i)}
    \mathbb{I}\{c_j=c_i\}
    +a
  }.
  \label{eq:ensemble-catboost-ordered-statistic}
\end{equation}
其中$p$是先验统计量，$a>0$控制平滑。样本自身的标签不会进入$z_i$；部署时则可使用训练集中的类别统计。低频类别仍会产生较大估计波动，平滑、多个排列和类别组合需要共同考虑。

图\ref{fig:ensemble-catboost-ordered-statistics}用一个排列位置说明这种隔离。计算当前类别$A$样本的编码时，只允许排列前缀中同为类别$A$的标签进入统计量；其他类别、当前标签和排列后缀都不参与这次编码。

\begin{figure}[htbp]
  \centering
  % CatBoost ordered target statistic for one sample in a random permutation.
% Schematic only; category symbols and labels are placeholders.
% Canvas: 112 mm x 58 mm. Dependencies: project palette and TikZ.
\begin{tikzpicture}[
  x=1mm,y=1mm,
  font=\figurefont\footnotesize,
  text=FigureInk,
  record/.style={
    rectangle,
    model input,
    minimum width=17mm,
    minimum height=14mm,
    text width=14mm,
    align=center,
    inner sep=1mm
  },
  current/.style={
    record,
    model training
  },
  future/.style={
    record,
    model context,
    text=FigureMuted
  },
  result/.style={
    rectangle,
    model output,
    minimum width=50mm,
    minimum height=12mm,
    text width=46mm,
    align=center,
    inner sep=1mm
  },
  flow/.style={
    model flow
  },
  use/.style={
    model flow
  }
]
  \path[use as bounding box] (0,0) rectangle (112,58);

  \draw[flow] (4,47)--(108,47)
    node[right,text=FigureInk] {$\pi$递增};

  \node[record] (p1) at (14,32)
    {$j_1$\\$c_{j_1}=A$\\$y_{j_1}$};
  \node[record] (p2) at (36,32)
    {$j_2$\\$c_{j_2}=B$\\$y_{j_2}$};
  \node[record] (p3) at (58,32)
    {$j_3$\\$c_{j_3}=A$\\$y_{j_3}$};
  \node[current] (pi) at (80,32)
    {当前样本$i$\\$c_i=A$};
  \node[future] (pf) at (102,32)
    {后续样本$j_f$\\$c_{j_f}=A$};
  \foreach \name in {p1,p2,p3,pi,pf} {
    \draw[FigureMuted,line width=0.7pt] (\name.north)--++(0,6);
  }

  \node[result] (statistic) at (47,9)
    {汇总同类前缀标签\\得到样本$i$的编码$z_i$};
  \draw[FigureInk,line width=.9pt] (p1.south)--(14,20);
  \draw[FigureInk,line width=.9pt] (p3.south)--(58,20);
  \draw[FigureInk,line width=.9pt] (14,20)--(58,20);
  \fill[FigureInk] (47,20) circle[radius=1.1pt];
  \draw[use] (47,20)--(statistic.north);
\end{tikzpicture}

  \caption{CatBoost有序目标统计的前缀约束。这里$\pi(j)$表示样本$j$在随机排列中的位置；计算当前样本$i$的编码时，只有满足$\pi(j)<\pi(i)$且$c_j=c_i$的标签通过实线汇入统计量。其他类别、当前标签和后续记录均不连接到统计量，表示它们在本次编码中被隔离。样本索引与类别符号仅用于说明机制。}
  \label{fig:ensemble-catboost-ordered-statistics}
\end{figure}

有序提升把相同的前缀思想用于梯度估计。直观地说，计算某个训练样本的当前预测与伪残差时，只使用排列中位于该样本之前的记录形成相应模型，减少“同一样本既参与拟合当前模型，又用于评价当前模型”的预测偏移。实际算法通过多个排列和前缀模型实现这一约束，训练状态与计算量也比普通单序列提升更复杂\scite{prokhorenkova2018catboost}

CatBoost常使用\term{对称树}（symmetric tree），即同一深度的所有节点采用相同分裂条件。深度为$d_{\mathrm{tree}}$的完整对称树至多有$2^{d_{\mathrm{tree}}}$片叶，路径可以编码为固定长度的二进制选择。这种结构便于批量预测，也限制了单棵树可表达的分区形状。类别编码、有序提升和对称树分别处理输入构造、序贯估计偏移和成员结构，不能把某一项的作用归给整个系统。

\subsection{模型对比与总结}
\label{sec:ensemble-boosting-comparison}

表\ref{tab:ensemble-boosting-comparison}先固定共同的序贯加法结构，再比较各方法修改哪一层。AdaBoost由加权分类错误生成样本权重；GBM由损失梯度生成伪残差；XGBoost对正则化树目标作二阶近似；LightGBM进一步压缩分裂统计的计算；CatBoost重点处理类别统计与序贯估计中的目标泄漏。后四者之间存在实现重叠，表中只列代表性机制，不把软件当前支持的全部选项写成永恒差异。

\begin{table*}[htbp]
  \centering
  \small
  \begin{tabularx}{\linewidth}{@{}>{\raggedright\arraybackslash}p{19mm}*{6}{>{\raggedright\arraybackslash}X}@{}}
    \toprule
    方法 & 本轮修正信号 & 典型成员 & 目标近似 & 树生长或分裂 & 类别特征重点 & 主要约束 \\
    \midrule
    AdaBoost & 加权错分 & 树桩或浅分类树 & 指数损失下的解析更新 & 由成员算法决定 & 需另行编码 & 对持续错分记录敏感 \\
    GBM／GBDT & 负梯度伪残差 & 回归树 & 一阶函数方向与叶内优化 & 通常贪心树生长 & 需随实现处理 & 学习率、轮数与树复杂度耦合 \\
    XGBoost & 一阶导与二阶导 & 正则化回归树 & 二阶局部目标 & 精确或分位数候选；稀疏默认方向 & 依实现与版本而定 & 轮次串行，近似与系统配置影响成本 \\
    LightGBM & 桶级梯度统计 & 直方图回归树 & 通常使用二阶统计 & 按叶生长；直方图候选 & 支持类别分组，机制依实现 & 分桶、GOSS与EFB引入近似条件 \\
    CatBoost & 有序统计与有序梯度 & 对称树 & 梯度或二阶统计 & 同层共享分裂条件 & 有序目标统计与组合 & 多排列和前缀状态增加训练成本 \\
    \bottomrule
  \end{tabularx}
  \caption{Boosting代表方法的机制比较。GBM是一般框架，GBDT指使用回归树成员的实例；XGBoost、LightGBM和CatBoost是树提升的具体模型与系统，不能与Boosting这一上位类别并列。}
  \label{tab:ensemble-boosting-comparison}
\end{table*}

% 跨栏表之后的扩展段需与三篇完整出处共同容纳，不能只按正文行数分页。
\Needspace{23\baselineskip}
随机梯度提升在每轮使用训练子样本，DART随机暂时移除部分已有树后计算修正，LogitBoost则从加法Logistic回归推导更新\scite{friedman2002stochastic,rashmi2015dart,friedman2000adaboost}理解这些扩展时，可以询问它改变了训练样本、已有成员、损失近似还是树结构。Boosting能够沿训练目标定向增加容量，却仍采用预先规定的加法组合形式；若已有成员来自不同模型族，并希望从它们的样本外表现中学习组合方式，就进入Stacking的范围。
